%% ── DantinoX paper code listings ────────────────────────────────────────────
%%
%% Include in the main document preamble:
%%   \usepackage{listings,xcolor}
%%   \input{code_listings_preamble}   % or copy the lstdefinestyle block below
%%
%% Then insert each figure with  %% ── DantinoX paper code listings ────────────────────────────────────────────
%%
%% Include in the main document preamble:
%%   \usepackage{listings,xcolor}
%%   \input{code_listings_preamble}   % or copy the lstdefinestyle block below
%%
%% Then insert each figure with  %% ── DantinoX paper code listings ────────────────────────────────────────────
%%
%% Include in the main document preamble:
%%   \usepackage{listings,xcolor}
%%   \input{code_listings_preamble}   % or copy the lstdefinestyle block below
%%
%% Then insert each figure with  %% ── DantinoX paper code listings ────────────────────────────────────────────
%%
%% Include in the main document preamble:
%%   \usepackage{listings,xcolor}
%%   \input{code_listings_preamble}   % or copy the lstdefinestyle block below
%%
%% Then insert each figure with  \input{code_listings}  or copy/paste.
%% ─────────────────────────────────────────────────────────────────────────────

%% ── Colour palette ───────────────────────────────────────────────────────────
\definecolor{lbg}{RGB}{245,247,252}
\definecolor{lkw}{RGB}{25,75,155}
\definecolor{lst}{RGB}{135,55,0}
\definecolor{lcm}{RGB}{90,90,90}
\definecolor{lnu}{RGB}{100,0,140}
\definecolor{lfr}{RGB}{190,200,220}

%% ── Shared listing style ─────────────────────────────────────────────────────
\lstdefinestyle{dxpy}{
  language        = Python,
  basicstyle      = \ttfamily\fontsize{7.0}{8.5}\selectfont,
  backgroundcolor = \color{lbg},
  commentstyle    = \color{lcm}\itshape,
  keywordstyle    = \color{lkw}\bfseries,
  stringstyle     = \color{lst},
  morekeywords    = {True,False,None},
  showstringspaces= false,
  breaklines      = true,
  breakatwhitespace=true,
  columns         = flexible,
  keepspaces      = true,
  frame           = single,
  framesep        = 2.5pt,
  rulecolor       = \color{lfr},
  xleftmargin     = 4pt,
  xrightmargin    = 4pt,
  aboveskip       = 3pt,
  belowskip       = 2pt,
}

%% ═════════════════════════════════════════════════════════════════════════════
%%  Listing 1 — Training
%% ═════════════════════════════════════════════════════════════════════════════
\begin{figure}[t]
\begin{lstlisting}[style=dxpy]
import dantinox as dx

# One field switches the paradigm; all other components are shared.
cfg = dx.ModelConfig(
    paradigm   = "discrete",    # "ar" | "discrete" | "continuous"
    attention  = "gqa", kv_heads   = 2,   # MHA | GQA | MLA
    ffn        = "mlp", use_swiglu = True, # mlp | moe
    dim        = 512,   n_heads    = 8,
    num_blocks = 12,    vocab_size = 32128,
    tp_size    = 2,               # tensor-parallel shards per replica
)
tcfg = dx.TrainingConfig(
    lr           = 3e-4,  epochs     = 10,
    batch_size   = 64,    grad_accum = 4,   # effective batch = 256
    optimizer    = "adamw",  lr_schedule  = "cosine",
    warmup_steps = 400,
    n_devices    = 4,     tp_size      = 2, # 2 DP x 2 TP = 4 GPUs
    dataset_name   = "wikitext",
    dataset_config = "wikitext-103-raw-v1",
)
run_dir = dx.Trainer(dx.Paradigm(cfg), tcfg).fit()
\end{lstlisting}
\caption{Training a discrete diffusion model with DantinoX. Switching to
  autoregressive or ELF requires only changing \texttt{paradigm=\char`\"ar\char`\"}
  or \texttt{\char`\"continuous\char`\"}, while the attention variant, FFN,
  multi-GPU configuration, and \texttt{Trainer} remain identical.}
\label{lst:training}
\end{figure}

%% ═════════════════════════════════════════════════════════════════════════════
%%  Listing 2 — Inference & Benchmarking
%% ═════════════════════════════════════════════════════════════════════════════
\begin{figure}[t]
\begin{lstlisting}[style=dxpy]
from dantinox import Generator

# Generator reads config.yaml from the run_dir and selects the
# correct inference algorithm (AR / Discrete / ELF) automatically.
gen_ar   = Generator("runs/ar_gqa_512d_12b",  seed=42)
gen_diff = Generator("runs/diff_gqa_512d_12b", seed=42)

# AR: token-by-token with static KV-cache (O(1) memory/step).
for tok in gen_ar.stream("HAMLET:\n", max_new_tokens=200,
                          top_k=40, temperature=0.8):
    print(tok, end="", flush=True)

# Discrete Diffusion: n_steps reverse-diffusion passes.
for chunk in gen_diff.stream("Language models will",
                              max_new_tokens=100, n_steps=50,
                              decoding_strategy="confidence"):
    print(chunk, end="", flush=True)

# Fast-dLLM: block-wise denoising with DualCache (1.4-2.1x speedup).
for chunk in gen_diff.stream("Language models will",
                              max_new_tokens=100, use_blocks=True):
    print(chunk, end="", flush=True)

# ── Performance analysis ──────────────────────────────────────────
# Analytical FLOPs breakdown (no GPU needed).
fl = dx.count_flops(cfg, seq_len=256)
print(f"total={fl.total/1e9:.2f}G  "
      f"attn={fl.attention/1e9:.2f}G  ffn={fl.ffn/1e9:.2f}G")

# Analytical + measured profiling: latency, throughput, MFU.
from flax import nnx
paradigm = dx.Paradigm(cfg)
model    = paradigm.build_model(nnx.Rngs(0))
report   = dx.profile(cfg, seq_len=256, batch_size=4,
                      n_warmup=5, n_runs=20, model=model)
print(report)          # latency_ms, tokens_per_sec, mfu

# Full benchmark suite: Throughput + Latency + Perplexity.
from dantinox.benchmarking import BenchmarkSuite, BenchmarkConfig
suite  = BenchmarkSuite.default(BenchmarkConfig(
    seq_lens=[128, 256, 512], batch_sizes=[1, 4, 16, 64],
    n_warmup=5, n_measure=20,
))
report = suite.run(paradigm, model, save_csv="results/bench.csv")
print(report.summary())
\end{lstlisting}
\caption{Inference and performance analysis with DantinoX. \texttt{Generator}
  provides a unified streaming API across all paradigms. \texttt{count\_flops}
  gives an analytical FLOPs breakdown; \texttt{profile} additionally measures
  GPU latency and MFU; \texttt{BenchmarkSuite.default()} sweeps throughput,
  latency, and perplexity across sequence lengths and batch sizes.}
\label{lst:inference}
\end{figure}

%% ═════════════════════════════════════════════════════════════════════════════
%%  Listing 3 — Reproducing a Cross-Paradigm Results Table
%% ═════════════════════════════════════════════════════════════════════════════
\begin{figure}[t]
\begin{lstlisting}[style=dxpy]
import dantinox as dx

GRID    = [(p, a) for p in ("ar", "discrete", "continuous")
                   for a in ("mha", "gqa", "mla")]
N_STEPS = 64   # matched inference budget across paradigms

rows = []
for paradigm, attn in GRID:
    for seed in (42, 43, 44):
        cfg  = dx.ModelConfig(paradigm=paradigm, attention=attn,
                              dim=512, n_heads=8, num_blocks=12)
        tcfg = dx.TrainingConfig(optimizer="muon", batch_size=256,
                                 dataset_name="wikitext", seed=seed)
        run_dir = dx.Trainer(dx.Paradigm(cfg), tcfg).fit(
            dataset_source="huggingface")

        gen     = dx.Generator(run_dir, seed=seed)
        samples = [gen.generate("", max_new_tokens=128, n_steps=N_STEPS)
                  for _ in range(100)]          # gen.seed varied per call
        rows.append(score(samples, paradigm))   # MAUVE, PPL, D-2, R-4, B-4

table = aggregate(rows, group=("paradigm", "attn"))  # mean +/- std, 3 seeds
write_latex(table, "tab_genquality_results.tex")
\end{lstlisting}
\vspace{-2mm}
\caption{\textsc{DantinoX} pipeline reproducing
  Table~\ref{tab:genquality}: one \texttt{ModelConfig}/\texttt{TrainingConfig}
  pair per cell of the paradigm~$\times$~attention grid, trained and scored
  through the same public API, with no paradigm-specific code.}
\label{lst:table2}
\end{figure}
  or copy/paste.
%% ─────────────────────────────────────────────────────────────────────────────

%% ── Colour palette ───────────────────────────────────────────────────────────
\definecolor{lbg}{RGB}{245,247,252}
\definecolor{lkw}{RGB}{25,75,155}
\definecolor{lst}{RGB}{135,55,0}
\definecolor{lcm}{RGB}{90,90,90}
\definecolor{lnu}{RGB}{100,0,140}
\definecolor{lfr}{RGB}{190,200,220}

%% ── Shared listing style ─────────────────────────────────────────────────────
\lstdefinestyle{dxpy}{
  language        = Python,
  basicstyle      = \ttfamily\fontsize{7.0}{8.5}\selectfont,
  backgroundcolor = \color{lbg},
  commentstyle    = \color{lcm}\itshape,
  keywordstyle    = \color{lkw}\bfseries,
  stringstyle     = \color{lst},
  morekeywords    = {True,False,None},
  showstringspaces= false,
  breaklines      = true,
  breakatwhitespace=true,
  columns         = flexible,
  keepspaces      = true,
  frame           = single,
  framesep        = 2.5pt,
  rulecolor       = \color{lfr},
  xleftmargin     = 4pt,
  xrightmargin    = 4pt,
  aboveskip       = 3pt,
  belowskip       = 2pt,
}

%% ═════════════════════════════════════════════════════════════════════════════
%%  Listing 1 — Training
%% ═════════════════════════════════════════════════════════════════════════════
\begin{figure}[t]
\begin{lstlisting}[style=dxpy]
import dantinox as dx

# One field switches the paradigm; all other components are shared.
cfg = dx.ModelConfig(
    paradigm   = "discrete",    # "ar" | "discrete" | "continuous"
    attention  = "gqa", kv_heads   = 2,   # MHA | GQA | MLA
    ffn        = "mlp", use_swiglu = True, # mlp | moe
    dim        = 512,   n_heads    = 8,
    num_blocks = 12,    vocab_size = 32128,
    tp_size    = 2,               # tensor-parallel shards per replica
)
tcfg = dx.TrainingConfig(
    lr           = 3e-4,  epochs     = 10,
    batch_size   = 64,    grad_accum = 4,   # effective batch = 256
    optimizer    = "adamw",  lr_schedule  = "cosine",
    warmup_steps = 400,
    n_devices    = 4,     tp_size      = 2, # 2 DP x 2 TP = 4 GPUs
    dataset_name   = "wikitext",
    dataset_config = "wikitext-103-raw-v1",
)
run_dir = dx.Trainer(dx.Paradigm(cfg), tcfg).fit()
\end{lstlisting}
\caption{Training a discrete diffusion model with DantinoX. Switching to
  autoregressive or ELF requires only changing \texttt{paradigm=\char`\"ar\char`\"}
  or \texttt{\char`\"continuous\char`\"}, while the attention variant, FFN,
  multi-GPU configuration, and \texttt{Trainer} remain identical.}
\label{lst:training}
\end{figure}

%% ═════════════════════════════════════════════════════════════════════════════
%%  Listing 2 — Inference & Benchmarking
%% ═════════════════════════════════════════════════════════════════════════════
\begin{figure}[t]
\begin{lstlisting}[style=dxpy]
from dantinox import Generator

# Generator reads config.yaml from the run_dir and selects the
# correct inference algorithm (AR / Discrete / ELF) automatically.
gen_ar   = Generator("runs/ar_gqa_512d_12b",  seed=42)
gen_diff = Generator("runs/diff_gqa_512d_12b", seed=42)

# AR: token-by-token with static KV-cache (O(1) memory/step).
for tok in gen_ar.stream("HAMLET:\n", max_new_tokens=200,
                          top_k=40, temperature=0.8):
    print(tok, end="", flush=True)

# Discrete Diffusion: n_steps reverse-diffusion passes.
for chunk in gen_diff.stream("Language models will",
                              max_new_tokens=100, n_steps=50,
                              decoding_strategy="confidence"):
    print(chunk, end="", flush=True)

# Fast-dLLM: block-wise denoising with DualCache (1.4-2.1x speedup).
for chunk in gen_diff.stream("Language models will",
                              max_new_tokens=100, use_blocks=True):
    print(chunk, end="", flush=True)

# ── Performance analysis ──────────────────────────────────────────
# Analytical FLOPs breakdown (no GPU needed).
fl = dx.count_flops(cfg, seq_len=256)
print(f"total={fl.total/1e9:.2f}G  "
      f"attn={fl.attention/1e9:.2f}G  ffn={fl.ffn/1e9:.2f}G")

# Analytical + measured profiling: latency, throughput, MFU.
from flax import nnx
paradigm = dx.Paradigm(cfg)
model    = paradigm.build_model(nnx.Rngs(0))
report   = dx.profile(cfg, seq_len=256, batch_size=4,
                      n_warmup=5, n_runs=20, model=model)
print(report)          # latency_ms, tokens_per_sec, mfu

# Full benchmark suite: Throughput + Latency + Perplexity.
from dantinox.benchmarking import BenchmarkSuite, BenchmarkConfig
suite  = BenchmarkSuite.default(BenchmarkConfig(
    seq_lens=[128, 256, 512], batch_sizes=[1, 4, 16, 64],
    n_warmup=5, n_measure=20,
))
report = suite.run(paradigm, model, save_csv="results/bench.csv")
print(report.summary())
\end{lstlisting}
\caption{Inference and performance analysis with DantinoX. \texttt{Generator}
  provides a unified streaming API across all paradigms. \texttt{count\_flops}
  gives an analytical FLOPs breakdown; \texttt{profile} additionally measures
  GPU latency and MFU; \texttt{BenchmarkSuite.default()} sweeps throughput,
  latency, and perplexity across sequence lengths and batch sizes.}
\label{lst:inference}
\end{figure}

%% ═════════════════════════════════════════════════════════════════════════════
%%  Listing 3 — Reproducing a Cross-Paradigm Results Table
%% ═════════════════════════════════════════════════════════════════════════════
\begin{figure}[t]
\begin{lstlisting}[style=dxpy]
import dantinox as dx

GRID    = [(p, a) for p in ("ar", "discrete", "continuous")
                   for a in ("mha", "gqa", "mla")]
N_STEPS = 64   # matched inference budget across paradigms

rows = []
for paradigm, attn in GRID:
    for seed in (42, 43, 44):
        cfg  = dx.ModelConfig(paradigm=paradigm, attention=attn,
                              dim=512, n_heads=8, num_blocks=12)
        tcfg = dx.TrainingConfig(optimizer="muon", batch_size=256,
                                 dataset_name="wikitext", seed=seed)
        run_dir = dx.Trainer(dx.Paradigm(cfg), tcfg).fit(
            dataset_source="huggingface")

        gen     = dx.Generator(run_dir, seed=seed)
        samples = [gen.generate("", max_new_tokens=128, n_steps=N_STEPS)
                  for _ in range(100)]          # gen.seed varied per call
        rows.append(score(samples, paradigm))   # MAUVE, PPL, D-2, R-4, B-4

table = aggregate(rows, group=("paradigm", "attn"))  # mean +/- std, 3 seeds
write_latex(table, "tab_genquality_results.tex")
\end{lstlisting}
\vspace{-2mm}
\caption{\textsc{DantinoX} pipeline reproducing
  Table~\ref{tab:genquality}: one \texttt{ModelConfig}/\texttt{TrainingConfig}
  pair per cell of the paradigm~$\times$~attention grid, trained and scored
  through the same public API, with no paradigm-specific code.}
\label{lst:table2}
\end{figure}
  or copy/paste.
%% ─────────────────────────────────────────────────────────────────────────────

%% ── Colour palette ───────────────────────────────────────────────────────────
\definecolor{lbg}{RGB}{245,247,252}
\definecolor{lkw}{RGB}{25,75,155}
\definecolor{lst}{RGB}{135,55,0}
\definecolor{lcm}{RGB}{90,90,90}
\definecolor{lnu}{RGB}{100,0,140}
\definecolor{lfr}{RGB}{190,200,220}

%% ── Shared listing style ─────────────────────────────────────────────────────
\lstdefinestyle{dxpy}{
  language        = Python,
  basicstyle      = \ttfamily\fontsize{7.0}{8.5}\selectfont,
  backgroundcolor = \color{lbg},
  commentstyle    = \color{lcm}\itshape,
  keywordstyle    = \color{lkw}\bfseries,
  stringstyle     = \color{lst},
  morekeywords    = {True,False,None},
  showstringspaces= false,
  breaklines      = true,
  breakatwhitespace=true,
  columns         = flexible,
  keepspaces      = true,
  frame           = single,
  framesep        = 2.5pt,
  rulecolor       = \color{lfr},
  xleftmargin     = 4pt,
  xrightmargin    = 4pt,
  aboveskip       = 3pt,
  belowskip       = 2pt,
}

%% ═════════════════════════════════════════════════════════════════════════════
%%  Listing 1 — Training
%% ═════════════════════════════════════════════════════════════════════════════
\begin{figure}[t]
\begin{lstlisting}[style=dxpy]
import dantinox as dx

# One field switches the paradigm; all other components are shared.
cfg = dx.ModelConfig(
    paradigm   = "discrete",    # "ar" | "discrete" | "continuous"
    attention  = "gqa", kv_heads   = 2,   # MHA | GQA | MLA
    ffn        = "mlp", use_swiglu = True, # mlp | moe
    dim        = 512,   n_heads    = 8,
    num_blocks = 12,    vocab_size = 32128,
    tp_size    = 2,               # tensor-parallel shards per replica
)
tcfg = dx.TrainingConfig(
    lr           = 3e-4,  epochs     = 10,
    batch_size   = 64,    grad_accum = 4,   # effective batch = 256
    optimizer    = "adamw",  lr_schedule  = "cosine",
    warmup_steps = 400,
    n_devices    = 4,     tp_size      = 2, # 2 DP x 2 TP = 4 GPUs
    dataset_name   = "wikitext",
    dataset_config = "wikitext-103-raw-v1",
)
run_dir = dx.Trainer(dx.Paradigm(cfg), tcfg).fit()
\end{lstlisting}
\caption{Training a discrete diffusion model with DantinoX. Switching to
  autoregressive or ELF requires only changing \texttt{paradigm=\char`\"ar\char`\"}
  or \texttt{\char`\"continuous\char`\"}, while the attention variant, FFN,
  multi-GPU configuration, and \texttt{Trainer} remain identical.}
\label{lst:training}
\end{figure}

%% ═════════════════════════════════════════════════════════════════════════════
%%  Listing 2 — Inference & Benchmarking
%% ═════════════════════════════════════════════════════════════════════════════
\begin{figure}[t]
\begin{lstlisting}[style=dxpy]
from dantinox import Generator

# Generator reads config.yaml from the run_dir and selects the
# correct inference algorithm (AR / Discrete / ELF) automatically.
gen_ar   = Generator("runs/ar_gqa_512d_12b",  seed=42)
gen_diff = Generator("runs/diff_gqa_512d_12b", seed=42)

# AR: token-by-token with static KV-cache (O(1) memory/step).
for tok in gen_ar.stream("HAMLET:\n", max_new_tokens=200,
                          top_k=40, temperature=0.8):
    print(tok, end="", flush=True)

# Discrete Diffusion: n_steps reverse-diffusion passes.
for chunk in gen_diff.stream("Language models will",
                              max_new_tokens=100, n_steps=50,
                              decoding_strategy="confidence"):
    print(chunk, end="", flush=True)

# Fast-dLLM: block-wise denoising with DualCache (1.4-2.1x speedup).
for chunk in gen_diff.stream("Language models will",
                              max_new_tokens=100, use_blocks=True):
    print(chunk, end="", flush=True)

# ── Performance analysis ──────────────────────────────────────────
# Analytical FLOPs breakdown (no GPU needed).
fl = dx.count_flops(cfg, seq_len=256)
print(f"total={fl.total/1e9:.2f}G  "
      f"attn={fl.attention/1e9:.2f}G  ffn={fl.ffn/1e9:.2f}G")

# Analytical + measured profiling: latency, throughput, MFU.
from flax import nnx
paradigm = dx.Paradigm(cfg)
model    = paradigm.build_model(nnx.Rngs(0))
report   = dx.profile(cfg, seq_len=256, batch_size=4,
                      n_warmup=5, n_runs=20, model=model)
print(report)          # latency_ms, tokens_per_sec, mfu

# Full benchmark suite: Throughput + Latency + Perplexity.
from dantinox.benchmarking import BenchmarkSuite, BenchmarkConfig
suite  = BenchmarkSuite.default(BenchmarkConfig(
    seq_lens=[128, 256, 512], batch_sizes=[1, 4, 16, 64],
    n_warmup=5, n_measure=20,
))
report = suite.run(paradigm, model, save_csv="results/bench.csv")
print(report.summary())
\end{lstlisting}
\caption{Inference and performance analysis with DantinoX. \texttt{Generator}
  provides a unified streaming API across all paradigms. \texttt{count\_flops}
  gives an analytical FLOPs breakdown; \texttt{profile} additionally measures
  GPU latency and MFU; \texttt{BenchmarkSuite.default()} sweeps throughput,
  latency, and perplexity across sequence lengths and batch sizes.}
\label{lst:inference}
\end{figure}

%% ═════════════════════════════════════════════════════════════════════════════
%%  Listing 3 — Reproducing a Cross-Paradigm Results Table
%% ═════════════════════════════════════════════════════════════════════════════
\begin{figure}[t]
\begin{lstlisting}[style=dxpy]
import dantinox as dx

GRID    = [(p, a) for p in ("ar", "discrete", "continuous")
                   for a in ("mha", "gqa", "mla")]
N_STEPS = 64   # matched inference budget across paradigms

rows = []
for paradigm, attn in GRID:
    for seed in (42, 43, 44):
        cfg  = dx.ModelConfig(paradigm=paradigm, attention=attn,
                              dim=512, n_heads=8, num_blocks=12)
        tcfg = dx.TrainingConfig(optimizer="muon", batch_size=256,
                                 dataset_name="wikitext", seed=seed)
        run_dir = dx.Trainer(dx.Paradigm(cfg), tcfg).fit(
            dataset_source="huggingface")

        gen     = dx.Generator(run_dir, seed=seed)
        samples = [gen.generate("", max_new_tokens=128, n_steps=N_STEPS)
                  for _ in range(100)]          # gen.seed varied per call
        rows.append(score(samples, paradigm))   # MAUVE, PPL, D-2, R-4, B-4

table = aggregate(rows, group=("paradigm", "attn"))  # mean +/- std, 3 seeds
write_latex(table, "tab_genquality_results.tex")
\end{lstlisting}
\vspace{-2mm}
\caption{\textsc{DantinoX} pipeline reproducing
  Table~\ref{tab:genquality}: one \texttt{ModelConfig}/\texttt{TrainingConfig}
  pair per cell of the paradigm~$\times$~attention grid, trained and scored
  through the same public API, with no paradigm-specific code.}
\label{lst:table2}
\end{figure}
  or copy/paste.
%% ─────────────────────────────────────────────────────────────────────────────

%% ── Colour palette ───────────────────────────────────────────────────────────
\definecolor{lbg}{RGB}{245,247,252}
\definecolor{lkw}{RGB}{25,75,155}
\definecolor{lst}{RGB}{135,55,0}
\definecolor{lcm}{RGB}{90,90,90}
\definecolor{lnu}{RGB}{100,0,140}
\definecolor{lfr}{RGB}{190,200,220}

%% ── Shared listing style ─────────────────────────────────────────────────────
\lstdefinestyle{dxpy}{
  language        = Python,
  basicstyle      = \ttfamily\fontsize{7.0}{8.5}\selectfont,
  backgroundcolor = \color{lbg},
  commentstyle    = \color{lcm}\itshape,
  keywordstyle    = \color{lkw}\bfseries,
  stringstyle     = \color{lst},
  morekeywords    = {True,False,None},
  showstringspaces= false,
  breaklines      = true,
  breakatwhitespace=true,
  columns         = flexible,
  keepspaces      = true,
  frame           = single,
  framesep        = 2.5pt,
  rulecolor       = \color{lfr},
  xleftmargin     = 4pt,
  xrightmargin    = 4pt,
  aboveskip       = 3pt,
  belowskip       = 2pt,
}

%% ═════════════════════════════════════════════════════════════════════════════
%%  Listing 1 — Training
%% ═════════════════════════════════════════════════════════════════════════════
\begin{figure}[t]
\begin{lstlisting}[style=dxpy]
import dantinox as dx

# One field switches the paradigm; all other components are shared.
cfg = dx.ModelConfig(
    paradigm   = "discrete",    # "ar" | "discrete" | "continuous"
    attention  = "gqa", kv_heads   = 2,   # MHA | GQA | MLA
    ffn        = "mlp", use_swiglu = True, # mlp | moe
    dim        = 512,   n_heads    = 8,
    num_blocks = 12,    vocab_size = 32128,
    tp_size    = 2,               # tensor-parallel shards per replica
)
tcfg = dx.TrainingConfig(
    lr           = 3e-4,  epochs     = 10,
    batch_size   = 64,    grad_accum = 4,   # effective batch = 256
    optimizer    = "adamw",  lr_schedule  = "cosine",
    warmup_steps = 400,
    n_devices    = 4,     tp_size      = 2, # 2 DP x 2 TP = 4 GPUs
    dataset_name   = "wikitext",
    dataset_config = "wikitext-103-raw-v1",
)
run_dir = dx.Trainer(dx.Paradigm(cfg), tcfg).fit()
\end{lstlisting}
\caption{Training a discrete diffusion model with DantinoX. Switching to
  autoregressive or ELF requires only changing \texttt{paradigm=\char`\"ar\char`\"}
  or \texttt{\char`\"continuous\char`\"}, while the attention variant, FFN,
  multi-GPU configuration, and \texttt{Trainer} remain identical.}
\label{lst:training}
\end{figure}

%% ═════════════════════════════════════════════════════════════════════════════
%%  Listing 2 — Inference & Benchmarking
%% ═════════════════════════════════════════════════════════════════════════════
\begin{figure}[t]
\begin{lstlisting}[style=dxpy]
from dantinox import Generator

# Generator reads config.yaml from the run_dir and selects the
# correct inference algorithm (AR / Discrete / ELF) automatically.
gen_ar   = Generator("runs/ar_gqa_512d_12b",  seed=42)
gen_diff = Generator("runs/diff_gqa_512d_12b", seed=42)

# AR: token-by-token with static KV-cache (O(1) memory/step).
for tok in gen_ar.stream("HAMLET:\n", max_new_tokens=200,
                          top_k=40, temperature=0.8):
    print(tok, end="", flush=True)

# Discrete Diffusion: n_steps reverse-diffusion passes.
for chunk in gen_diff.stream("Language models will",
                              max_new_tokens=100, n_steps=50,
                              decoding_strategy="confidence"):
    print(chunk, end="", flush=True)

# Fast-dLLM: block-wise denoising with DualCache (1.4-2.1x speedup).
for chunk in gen_diff.stream("Language models will",
                              max_new_tokens=100, use_blocks=True):
    print(chunk, end="", flush=True)

# ── Performance analysis ──────────────────────────────────────────
# Analytical FLOPs breakdown (no GPU needed).
fl = dx.count_flops(cfg, seq_len=256)
print(f"total={fl.total/1e9:.2f}G  "
      f"attn={fl.attention/1e9:.2f}G  ffn={fl.ffn/1e9:.2f}G")

# Analytical + measured profiling: latency, throughput, MFU.
from flax import nnx
paradigm = dx.Paradigm(cfg)
model    = paradigm.build_model(nnx.Rngs(0))
report   = dx.profile(cfg, seq_len=256, batch_size=4,
                      n_warmup=5, n_runs=20, model=model)
print(report)          # latency_ms, tokens_per_sec, mfu

# Full benchmark suite: Throughput + Latency + Perplexity.
from dantinox.benchmarking import BenchmarkSuite, BenchmarkConfig
suite  = BenchmarkSuite.default(BenchmarkConfig(
    seq_lens=[128, 256, 512], batch_sizes=[1, 4, 16, 64],
    n_warmup=5, n_measure=20,
))
report = suite.run(paradigm, model, save_csv="results/bench.csv")
print(report.summary())
\end{lstlisting}
\caption{Inference and performance analysis with DantinoX. \texttt{Generator}
  provides a unified streaming API across all paradigms. \texttt{count\_flops}
  gives an analytical FLOPs breakdown; \texttt{profile} additionally measures
  GPU latency and MFU; \texttt{BenchmarkSuite.default()} sweeps throughput,
  latency, and perplexity across sequence lengths and batch sizes.}
\label{lst:inference}
\end{figure}

%% ═════════════════════════════════════════════════════════════════════════════
%%  Listing 3 — Reproducing a Cross-Paradigm Results Table
%% ═════════════════════════════════════════════════════════════════════════════
\begin{figure}[t]
\begin{lstlisting}[style=dxpy]
import dantinox as dx

GRID    = [(p, a) for p in ("ar", "discrete", "continuous")
                   for a in ("mha", "gqa", "mla")]
N_STEPS = 64   # matched inference budget across paradigms

rows = []
for paradigm, attn in GRID:
    for seed in (42, 43, 44):
        cfg  = dx.ModelConfig(paradigm=paradigm, attention=attn,
                              dim=512, n_heads=8, num_blocks=12)
        tcfg = dx.TrainingConfig(optimizer="muon", batch_size=256,
                                 dataset_name="wikitext", seed=seed)
        run_dir = dx.Trainer(dx.Paradigm(cfg), tcfg).fit(
            dataset_source="huggingface")

        gen     = dx.Generator(run_dir, seed=seed)
        samples = [gen.generate("", max_new_tokens=128, n_steps=N_STEPS)
                  for _ in range(100)]          # gen.seed varied per call
        rows.append(score(samples, paradigm))   # MAUVE, PPL, D-2, R-4, B-4

table = aggregate(rows, group=("paradigm", "attn"))  # mean +/- std, 3 seeds
write_latex(table, "tab_genquality_results.tex")
\end{lstlisting}
\vspace{-2mm}
\caption{\textsc{DantinoX} pipeline reproducing
  Table~\ref{tab:genquality}: one \texttt{ModelConfig}/\texttt{TrainingConfig}
  pair per cell of the paradigm~$\times$~attention grid, trained and scored
  through the same public API, with no paradigm-specific code.}
\label{lst:table2}
\end{figure}
